\section{A depth-independent lower bound on every critical rate}
\label{sec:critical-universal-baseline}

At every prior scale $0<t=t(d,L,N)\le1$, the global clean-margin event has an
exponentially small probability, uniformly in depth. The exponent has a
universal positive lower bound as the input dimension grows. This result
does not compare the excess rates of different Boolean functions and is
not a posterior sample-complexity lower bound.

Throughout this section the hidden activation is $\tanh$, the input is
$\bar x=(1,x)$ for $x\in\{-1,1\}^d$, the first-layer variance is $t^2/d$,
the subsequent hidden-weight variance is $t^2/N$, and the readout variance
is $t^2/N^2$. The bounds hold at each specified $(d,L,N)$, so $t$
may follow any prescribed schedule in these parameters; it need not remain
constant as width varies. At fixed $d,L$, write this schedule as $t(N)$.
The critical model of the rest of the paper is the constant schedule
$t(N)=1$. There are $L-1$ hidden layers, with $L\ge2$.
Write $A_{N,f}$ for the event $f(x)z(x)>1/2$ at every cube input, and
$A_N=\bigcup_f A_{N,f}$ for the global clean-margin event.

\begin{lemma}[A scalar exponential potential]
\label{lem:critical-exponential-potential}
For $0<t\le1$, $0\le a\le1$, and $G\sim\mathcal N(0,t^2q)$, $q\ge0$,
\[
 \mathbb E\exp\{(2a/t^2)\tanh^2G\}\le e^{2aq/t^2}.
\]
If $(G_x,G_y)$ is any centered Gaussian pair, $s\in\{-1,1\}$,
and $t^2v=\operatorname{Var}[(G_x+sG_y)/2]$, then
\[
 \mathbb E\exp\left\{\frac{2a}{t^2}\left(
       \frac{\tanh G_x+s\tanh G_y}{2}\right)^2\right\}
 \le e^{2av/t^2}.
\]
\end{lemma}

\begin{proof}
The case $a=0$ is immediate. Otherwise put $\lambda=2a/t^2$ and
$F_\lambda(z)=\exp(\lambda\tanh^2z)$. Writing
$u=\tanh^2z\in[0,1)$, direct differentiation gives
\[
 \frac{F_\lambda''(z)}{F_\lambda(z)}
 =2\lambda(1-u)[1-3u+2\lambda u(1-u)]\le\frac{2\lambda}{t^2}.
\]
To check the inequality, multiply its bracket by $t^2$. The resulting
quantity $t^2(1-u)(1-3u)+4a u(1-u)^2$ is affine in $t^2\in[0,1]$.
At the two endpoints it is at most $16/27$ and
$1-5u^2+4u^3\le1$, respectively.
The Gaussian heat equation and Gronwall's inequality therefore give
$\mathbb E F_\lambda(\sqrt r Z)\le\exp(\lambda r/t^2)$ for $r\ge0$.
All derivatives needed here are bounded for each fixed $\lambda$.
Taking $r=t^2q$ proves the first assertion.
The elementary identities for the sum and difference of two hyperbolic
tangents imply
\[
 \left|\frac{\tanh a+s\tanh b}{2}\right|
 \le\tanh\left|\frac{a+sb}{2}\right|.
\]
For example, after writing $a=A+D$, $b=A-D$, the half-sum equals
$\sinh(2A)/[\cosh(2A)+\cosh(2D)]$, whose absolute value is at most
$\tanh|A|$. The half-difference is identical with $A,D$ interchanged.
Applying the scalar assertion to $(G_x+sG_y)/2$ proves the claim,
without any nonsingularity assumption.
\end{proof}

\begin{theorem}[Finite-width clean-margin upper bounds for all $\gamma\ge1$]
\label{thm:critical-clean-upper}
For every $d\ge1$, $N\ge1$, $L\ge2$, $0<t\le1$, and every Boolean function $f$,
\[
 \mu_{N,t}(A_{N,f})\le
       \exp\{-(N/t^2)\psi(d)\},
 \qquad
 \psi(d)=\begin{cases}d/8,&d\le4,\\1-2/d,&d\ge4.\end{cases}
\]
The unconditioned Gaussian probability of global clean margin satisfies
\[
 \mu_{N,t}(A_N)\le
       \min\{1,\ 2^{d+1}\exp[-(N/t^2)\psi(d)]\}.
\]
In the usual scaling $t^2=N^{1-\gamma}$, the exponential speed $N/t^2$
is $N^\gamma$, and the assumption is $\gamma\ge1$.
The first bound holds with no factor depending on depth or width outside
its exponential. The simpler exponent $1-2/d$ is valid in every dimension.
\end{theorem}

\begin{proof}
Fix two inputs $x,y$, a choice $s\in\{-1,1\}$, and $a\in(0,1]$.
For the empirical
hidden feature covariance $Q_j$, set
\[
 v_j=\frac14\bigl(Q_j(x,x)+Q_j(y,y)+2sQ_j(x,y)\bigr),
 \qquad 0\le j\le L-1,
\]
where $Q_0(x,y)=\bar x\cdot\bar y/d$. Conditional on the previous
layer, the $N$ preactivation pairs are independent centered Gaussians
with covariance given by $t^2Q_{j-1}$. Lemma~\ref{lem:critical-exponential-potential}
therefore gives
\[
 \mathbb E[\exp((2aN/t^2)v_j)\mid\text{previous layers}]
       \le\exp((2aN/t^2)v_{j-1}).
\]
Consequently $\mathbb E\exp((2aN/t^2)v_{L-1})\le\exp((2aN/t^2)v_0)$.

Conditional on the hidden layers, the readout combination
$Z=(z(x)+sz(y))/2$ is centered Gaussian with variance $t^2v_{L-1}/N$.
For either specified sign of $Z$, its probability of exceeding $1/2$
in that sign is at most $\exp[-N/(8t^2v_{L-1})]$, interpreted as zero
when $v_{L-1}=0$. Since
\[
 \frac1{8v}\ge\sqrt a-2av\qquad(v>0),
\]
the unconditional one-sided probability is at most
$\exp[-(N/t^2)(\sqrt a-2av_0)]$.

If $f$ is nonconstant, some adjacent cube vertices have opposite
labels. Choose such $x,y$ and $s=-1$. Their input difference has
squared norm four, so $v_0=1/d$, and $A_{N,f}$ implies a specified
sign of $Z$ exceeds $1/2$. If $f$ is constant, use antipodal inputs
$y=-x$ and $s=+1$. Their augmented-input half-sum is the bias
coordinate, again giving $v_0=1/d$. Maximizing $\sqrt a-2a/d$ over
$a\in[0,1]$ gives $\psi(d)$, proving the first assertion.

For the second assertion fix a spanning tree of the cube graph and
one antipodal pair. Every nonconstant sign function changes sign on
one of the tree's $2^d-1$ edges; a constant function has equal signs
on the fixed antipodal pair. Each case implies $|Z|>1/2$ for its
corresponding difference or sum. The two-sided bound and a union bound
over these $2^d$ pairs give the result.
\end{proof}

\begin{corollary}[Universal lower critical-rate baseline]
\label{cor:critical-rate-baseline}
For the critical variational rate functional $J_L$ with output margin
$1/2$, every Boolean function satisfies
\[
 J_L(f)\ge\psi(d)\ge1-\frac2d\qquad(L\ge2).
\]
\end{corollary}

\begin{proof}
This can be read directly from the finite-width upper bound whenever
the corresponding large-deviation formula is available. There is also
a direct proof from the variational functional, which avoids any
interchange of width and depth limits. Choose the signed pair from the
preceding proof, put $b=(e_x+se_y)/2$, and write $v_j=b^TQ_jb$ for
the covariances along any admissible variational path.
For any $a\in(0,1]$, the entropy variational inequality and
Lemma~\ref{lem:critical-exponential-potential} at $t=1$ give, for every hidden layer,
\[
 D(\nu_j\Vert\mathsf M_{Q_{j-1}})\ge2a(v_j-v_{j-1}).
\]
The output constraints imply $|b^Tz|\ge1/2$. If the readout cost is
finite, Cauchy--Schwarz on the range of $Q_{L-1}$ gives
\[
 \tfrac12 z^TQ_{L-1}^{+}z\ge\frac1{8v_{L-1}}.
\]
The inequality is immediate if this cost is infinite. Summing the
entropy inequalities and minimizing over $v=v_{L-1}>0$ now yields
\[
 J_L(f)\ge\inf_{v>0}\left(2av+\frac1{8v}\right)-2av_0
             =\sqrt a-\frac{2a}{d}.
\]
Maximizing in $a$ gives $\psi(d)$.
\end{proof}

For completeness, the same argument with an arbitrary required output
margin $r>0$ replaces $\psi(d)$ by
\[
 \psi_r(d)=\max_{0\le u\le1}(2ru-2u^2/d)
 =\begin{cases}r^2d/2,&rd\le2,\\2r-2/d,&rd\ge2.\end{cases}
\]
This applies to both finite-width bounds and the corresponding critical
variational rate. No change to the network architecture is needed.

The lower bound supplies a common baseline for rate comparisons. It
does not identify the minimum rate, bound the difference between a
target's rate and a competitor's rate, or show that the posterior is
diffuse among functions of bounded rate. In particular, even a small
positive rate gap is amplified by width in the critical posterior.
